\section{Approach}\label{sec:method}
% [M1-M11 L1-L4 CT1-CT3 R9 R10 W6 SC1 DL4] Internal stage names are explained in Appendix~\ref{app:glossary}.
Figure~\ref{fig:overview} shows the pipeline. We describe the model and the two training steps, then the variants used as controls and extensions.

\begin{figure*}[t]
\centering
\begin{tikzpicture}[font=\small, node distance=5mm, box/.style={draw, rounded corners=2pt, align=center, minimum height=9mm, inner sep=3pt}, arr/.style={-{Stealth[length=2mm]}, thick}, lab/.style={font=\scriptsize, align=center}]
\node[box] (g) {Gemma 3\\ 270M};
\node[box, right=of g] (grow) {grown to\\ 385M parameters};
\node[box, right=of grow, fill=gray!15] (sft) {\textbf{SFT baseline}\\ 125 language tags};
\node[box, right=of sft, fill=blue!12] (rl) {\textbf{RL model}\\ quality-score rewards};
\node[box, right=of rl] (sci) {science and\\ long-document\\ variants};
\node[box, below=9mm of sft, fill=orange!15] (ctl) {\textbf{SFT-only control}\\ same GPU time as RL};
\node[box, above=9mm of sft] (dec) {decoding tweaks\\ (repetition penalty, beam, MBR)};
\draw[arr] (g) -- (grow); \draw[arr] (grow) -- (sft); \draw[arr] (sft) -- (rl); \draw[arr] (rl) -- (sci);
\draw[arr] (sft) -- (ctl); \draw[arr, dashed] (dec) -- (sft);
\node[lab, below=1mm of ctl] {RQ3: compare with the RL model};
\node[lab, above=1mm of dec] {RQ2: apply to the SFT baseline and the RL model};
\node[lab, above=1.5mm of rl] {RQ1, RQ4};
\end{tikzpicture}
\caption{Overview. A 270M model is grown to 385M parameters and fine-tuned (the SFT baseline), then trained with reinforcement learning (the RL model). Two controls ask what else could give the same gain: decoding changes (RQ2) and more fine-tuning for the same GPU time (RQ3).}\label{fig:overview}
\end{figure*}

\textbf{Model.} We start from Gemma-3-270m-it \citep{gemmateam2025gemma3} and add layers to reach 385M parameters (39 layers) \citep[cf.][]{wu2024llamapro,kim2023solar}. The task is English to one of 125 languages. A single extra token in the vocabulary names the target language, and the model answers with that token followed by the translation.

\textbf{Step 1: supervised fine-tuning (the SFT baseline).} We train on English-to-African sentence pairs and instruction data. Before training, we remove every pair that matches the text of any benchmark we later use, checked on both sides of the pair. Part of the data comes from NLLB-mined sentences and a collection of translated instructions; we report their share and licences in the Ethical considerations. Training uses the Muon optimiser \citep{liu2025muon} and takes about 45 hours on one A100 GPU (Appendix~\ref{app:hyper}).

\textbf{Step 2: reinforcement learning (the RL model).} We use GRPO \citep{shao2024deepseekmath}. For each English text the model samples 16 translations. Each translation receives a reward, and the model is nudged towards translations that score above the average of their group and away from those below it; a penalty keeps the model close to its starting point. Three stages follow one another, and the first does most of the work (Section~\ref{sec:rq4}). The reward has three parts. (i) A translation that fails a hard check scores the minimum: wrong or missing language tag, no end of text, a copy of the input, repetition, the wrong language, or an implausible length. (ii) Otherwise a quality score without a reference translation rates the output. We combine two quality models \citep{wang2024africomet,li2025ssacomet}, taking their mean unless they disagree strongly, and for languages without a validated quality model we reward agreement with the group's other samples. (iii) Small penalties discourage, among other things, unusual length, missing numbers or names, and missing diacritics. Exact settings are in Appendix~\ref{app:hyper}.

\textbf{Variants.} The \emph{SFT-only control} continues SFT from the baseline for the same GPU time as the RL stages, using half replay data and half extra English-to-X pairs for 19 languages \citep{gritta2026translatepsy}. The \emph{science variant} adds a short fine-tuning step on science text from AfriScience-MT \citep{abdulmumin2026afriscience} and then a second RL run; the \emph{long-document variant} adds RL on long inputs. Internal names used in the appendix tables are listed in Appendix~\ref{app:glossary}.
